\chapter{PROJECT REQUIREMENT SPECIFICATIONS}

\section{Hardware Requirements}
\begin{table}[H]
\centering
\caption{Hardware Requirements}
\begin{tabular}{|p{4cm}|p{9cm}|}
\hline
\textbf{Component} & \textbf{Recommended Specification} \\
\hline
Processor & Multi-core CPU such as Intel i5/i7, AMD Ryzen, or Google Colab CPU runtime. \\
\hline
Memory & Minimum 8 GB RAM; 12 GB or higher recommended for the full Kaggle credit card fraud dataset. \\
\hline
Storage & Minimum 2 GB free space for dataset, checkpoints, plots, and report artifacts. \\
\hline
GPU & Optional. The MLP can use CUDA if available, but the project is designed to run on CPU. \\
\hline
Network & Required only for installing dependencies or using cloud notebook environments. Training itself is simulated locally. \\
\hline
\end{tabular}
\end{table}

\section{Software Requirements}
\begin{table}[H]
\centering
\caption{Software Requirements}
\begin{tabular}{|p{4cm}|p{9cm}|}
\hline
\textbf{Software} & \textbf{Purpose} \\
\hline
Python 3.10 or higher & Main implementation language. \\
\hline
Flower & Federated Learning client-server simulation and FedAvg strategy. \\
\hline
PyTorch & MLP implementation and neural network training. \\
\hline
scikit-learn & Logistic Regression, preprocessing, train-test split, metrics, and model utilities. \\
\hline
XGBoost & Client-local gradient boosted tree models. \\
\hline
pandas and NumPy & Dataset loading, tabular processing, numerical arrays. \\
\hline
matplotlib & ROC, PR, confusion matrix, and feature importance plots. \\
\hline
FastAPI and Uvicorn & Lightweight inference API. \\
\hline
SHAP & XGBoost explanation summaries. \\
\hline
\end{tabular}
\end{table}

\section{Functional Requirements}
\begin{enumerate}
    \item The system shall load CSV fraud datasets from the \texttt{data/} directory.
    \item The system shall automatically detect or validate the fraud label column.
    \item The system shall preprocess transaction features using imputation, scaling, and categorical encoding.
    \item The system shall partition the training data into multiple simulated federated clients.
    \item Each client shall train Logistic Regression and MLP models locally.
    \item The server shall aggregate Logistic Regression and MLP parameters using weighted FedAvg.
    \item Each client shall train a local XGBoost model for tabular fraud detection.
    \item The final predictor shall combine Logistic Regression, MLP, and XGBoost probability outputs using configurable ensemble weights.
    \item The system shall compute accuracy, precision, recall, F1-score, ROC-AUC, PR-AUC, and confusion matrix.
    \item The system shall save trained models, preprocessing pipeline, ensemble configuration, metrics, and plots.
    \item The API shall expose a \texttt{/predict} endpoint for real-time fraud scoring after training.
\end{enumerate}

\section{Non-Functional Requirements}
\begin{itemize}
    \item \textbf{Privacy:} Raw transaction data must remain inside client partitions during federated training.
    \item \textbf{Reproducibility:} Configuration values and random seeds must control training behaviour.
    \item \textbf{Modularity:} Data loading, models, clients, server, evaluation, and inference must be separated into reusable modules.
    \item \textbf{Scalability:} Number of clients, rounds, local epochs, and ensemble weights must be configurable.
    \item \textbf{Maintainability:} Code should use classes, type hints, checkpoint files, and clean project structure.
    \item \textbf{Interpretability:} The system should provide feature importance, SHAP summaries, and model contribution scores.
\end{itemize}



